\begin{frame}{결정경계: $k$를 키우면 어떻게 변하나}
  \begin{columns}[T]
    \begin{column}{0.55\textwidth}
      \centering
      \includegraphics[width=\linewidth]{ch04_knn_decision_boundary.png}
    \end{column}
    \begin{column}{0.43\textwidth}
      \begin{itemize}\small
        \item 두 원형 blob(클래스 0: 왼쪽 위, 1: 오른쪽 아래) 200개
          점, $k=1,5,40$의 결정경계
        \item \kb{$k=1$}: 훈련점 주변에 \textbf{톱니바퀴}처럼
          복잡하게 접힘 --- 각 점의 ``영향권''(보로노이 셀) 경계가
          그대로 결정경계
        \item \kb{$k=5$}: 톱니가 눈에 띄게 \textbf{매끄러워짐}
          (한두 점의 노이즈는 5표에서 다수결을 못 뒤집음)
        \item \kb{$k=40$}(훈련점 140개 중 1/3): \textbf{거의
          직선} $\to$ 지역 패턴 무시, 과소적합의 공간적 모습
      \end{itemize}
    \end{column}
  \end{columns}
  \vskip0.6em
  매끄러움이 어느 시점부터는 \textbf{정보를 잃는 매끄러움}
  (진짜 패턴 지우기)이 되는지 = \textbf{검증 정확도}로 판단
  (4.2절에서 실제로 $k$ 튜닝).
\end{frame}
